\section{系统模型与问题建模}
\label{rq3:sec:model}

\subsection{网络与语义任务模型}
\label{rq3:subsec:tasks}

考虑由 $N$ 个感知节点与一个接收中心组成的协同通信系统。感知节点集合记为 $\mathcal{N}=\{1,\ldots,N\}$，接收中心记为 $d$。系统依次处理 $N_{\mathrm{task}}$ 个任务，任务索引为 $k=1,\ldots,N_{\mathrm{task}}$。每个任务具有相对截止时间 $D_k$，其通信过程结束后，系统转入下一个任务。节点发送能量在任务之间连续扣减，接收状态和确认记录则在每个任务开始时重新初始化。

对于任务 $k$，上层候选选择模块根据局部观测与语义粒度给出候选证据集合
\begin{equation}
  \mathcal{M}_k=\{m_{k,1},m_{k,2},\ldots,m_{k,K_k}\}.
  \label{rq3:eq:candidates}
\end{equation}
每条证据 $m$ 包含发送节点 $s(m)$、长度 $\ell_m$、语义粒度和接收端信念更新量。不同粒度对应标签、简洁声明、摘要和完整证据。本章固定候选内容、候选次序和融合规则，控制器只调整这些证据的发送过程。选择来自另一节点的候选，即对应在给定候选集合内切换发送节点。

设接收中心实际已经获得的唯一证据集合为 $S_{k,t}\subseteq\mathcal{M}_k$，其中 $t$ 表示当前任务内的决策时刻。接收状态空间为
\begin{equation}
  \states_k=2^{\mathcal{M}_k},\qquad |\states_k|=2^{K_k}.
  \label{rq3:eq:state-space}
\end{equation}
接收中心按固定候选次序融合 $S_{k,t}$ 中的证据，并输出类别 $c_k(S_{k,t})$。相同证据无论成功接收多少次，都只参与一次融合。这一设定区分了通信层的重复发送与语义层的新增信息，使最终分类由证据集合决定，不随同一集合的到达次序变化。

记任务参考类别为 $y_k$，集合 $S$ 对应的真实任务效用定义为
\begin{equation}
  u_k(S)=\ind\{S\ne\varnothing,\ c_k(S)=y_k\}.
  \label{rq3:eq:true-utility}
\end{equation}
空接收集合的效用设为零，表示任务没有通过本次通信获得有效服务。$y_k$ 只用于离线效用学习和实验评价，不属于在线控制器的观测。由于在线无法直接计算式~\eqref{rq3:eq:true-utility}，第~\ref{rq3:subsec:utility} 节通过学习模型估计各接收集合的效用。

本章考虑短任务内持续存在的单跳接触。数据动作串行执行，反馈使用独立控制信道；不同尝试之间的速率扰动相互独立。候选描述及条件融合效果作为预先给定的调度信息，用于计算假设接收集合的语义特征。实验单独计入数据与确认的发送能耗，候选元数据分发、感知推理和飞行能耗不在控制范围内。这些假设使评价集中于给定语义内容后的交付决策。

\subsection{分组传输与随机送达模型}
\label{rq3:subsec:channel}

设候选 $m$ 的有效预测速率为 $\bar r_{k,m}$。链路预测由节点距离、有效带宽和发射功率给出。在本章使用的视距链路近似下，
\begin{equation}
  h_{k,i}=\frac{\rho_0}{\max\{d_{k,i},1\}^2},\qquad
  \bar r_{k,i}=W_k\log_2\left(1+\frac{P^{\mathrm d}h_{k,i}}{N_0W_k}\right),
  \label{rq3:eq:rate}
\end{equation}
其中 $d_{k,i}$ 为节点 $i$ 到中心的距离，$\rho_0$ 为参考信道增益，$N_0$ 为噪声功率谱密度，$P^{\mathrm d}$ 为数据发射功率，$W_k$ 表示任务对应的有效带宽。短任务内以该预测值计算动作时长，跨任务的距离与可用带宽变化则反映在不同的预测速率中。

语义长度以 token 数表示。设平均每个 token 占用 $b$ 字节，分组有效载荷上限为 $L_{\mathrm{MTU}}$ 字节，分组头部长度为 $H$ 字节，则实际发送比特数为
\begin{equation}
  B_m=8\left(b\ell_m+H\left\lceil\frac{b\ell_m}{L_{\mathrm{MTU}}}\right\rceil\right).
  \label{rq3:eq:bits}
\end{equation}
该换算同时保留了消息长度与协议头部开销。一次发送使用的时长为
\begin{equation}
  \tau_{k,m}=\min\left\{\tau_{\max},\ \kappa\frac{B_m}{\bar r_{k,m}}\right\},
  \label{rq3:eq:duration}
\end{equation}
其中 $\kappa$ 是发送时长裕量，$\tau_{\max}$ 是单次占用时长上限。控制器只允许完整时长能够容纳于剩余任务窗口内的发送动作。

为表示每次尝试的速率波动，令实际速率满足
\begin{equation}
  r_{k,m}^{(a)}=\bar r_{k,m}\exp\left(\sigma_r Z_a-\frac{\sigma_r^2}{2}\right),
  \qquad Z_a\sim\mathcal{N}(0,1),
  \label{rq3:eq:random-rate}
\end{equation}
其中 $a$ 为尝试编号，$\sigma_r$ 控制波动强度。若在预留时长内能够传完全部比特，则本次证据送达。相应成功概率为
\begin{equation}
  p_{k,m}=\Prb\{r_{k,m}^{(a)}\tau_{k,m}\ge B_m\}
  =\Phi\left(\frac{\ln(\bar r_{k,m}\tau_{k,m}/B_m)-\sigma_r^2/2}{\sigma_r}\right),
  \label{rq3:eq:success-probability}
\end{equation}
其中 $\Phi(\cdot)$ 为标准正态分布函数；当 $\sigma_r=0$ 时，上式退化为确定的容量判定。实验直接依据式~\eqref{rq3:eq:success-probability} 进行分组成功抽样。相同证据的不同发送尝试独立抽样，失败尝试消耗完整的预留发送时长。

本模型将未完整送达的证据视为擦除，不保存失败传输中的软信息。因此，重传的作用是提高证据至少成功送达一次的概率，而接收端不会对多个失败波形执行 Chase 合并或增量冗余译码。

\subsection{累计确认与可观测历史}
\label{rq3:subsec:feedback}

每次发送尝试结束后，接收中心生成一个累计确认帧。协议假设短控制头能够标识本次调度的发送尝试，因此即使载荷失败，也可返回包含当前接收状态的累计位图。反馈帧携带尝试序号 $j$ 及该时刻的接收集合快照 $S^{\mathrm{ack}}_j$，既可确认已收到的证据，也可表明某条已经尝试的证据仍未送达。

记第 $j$ 次数据发送结束时间为 $t_j^{\mathrm{end}}$，反馈帧长度为 $B^{\mathrm f}$ 比特。反馈按对应链路的预测速率串行发送，完成时间满足
\begin{equation}
  v_j=\max\{t_j^{\mathrm{end}},v_{j-1}\}+\tau_j^{\mathrm f},
  \qquad \tau_j^{\mathrm f}=\frac{B^{\mathrm f}}{\bar r_{k,m_j}},
  \label{rq3:eq:feedback-serialization}
\end{equation}
其中 $v_{-1}=0$。在串行发送之后，反馈以概率 $\epsilon_{\mathrm f}$ 丢失；若未丢失，则经过额外时延 $L_j$ 到达发送端，
\begin{equation}
  L_j\sim\mathcal{U}\left(\frac{\bar d_{\mathrm f}}{2},\frac{3\bar d_{\mathrm f}}{2}\right),
  \qquad t_j^{\mathrm{ack}}=v_j+L_j.
  \label{rq3:eq:feedback-delay}
\end{equation}
数据成功、反馈丢失和反馈额外时延分别抽样。累计确认可能乱序到达，控制器保留序号最大的有效快照，较旧位图不覆盖较新信息。

在决策时刻 $t$，可观测历史 $\history_{k,t}$ 包含候选描述、链路预测、已执行发送记录、已经到达的确认、剩余时间和能量状态。尚未到达的位图内容以及真实接收集合不属于该历史。记
\begin{equation}
  b_{k,t}(S)=\Prb\{S_{k,t}=S\mid\history_{k,t}\},\qquad S\in\states_k,
  \label{rq3:eq:belief-definition}
\end{equation}
为控制器对接收状态的信念。由于反馈过程与数据是否成功独立，确认尚未到达不能直接视为否定确认；它主要改变未来等待能够获得信息的概率。

\subsection{能量状态与动作约束}
\label{rq3:subsec:energy}

节点 $i$ 的初始发送能量为 $E_i^{(0)}$。一次数据发送及其反馈所消耗的能量分别为
\begin{equation}
  e_{k,m}^{\mathrm d}=P^{\mathrm d}\tau_{k,m},\qquad
  e_{k,m}^{\mathrm f}=P^{\mathrm f}\tau_{k,m}^{\mathrm f}.
  \label{rq3:eq:transmit-energy}
\end{equation}
数据发送能量从对应感知节点的剩余电量中扣除。接收中心采用独立供能，反馈发送能量计入系统通信成本。已经生成的反馈帧即使丢失或晚于任务结束才到达，也计入反馈能耗，以反映协议实际发起的发送开销。

设第 $k$ 个任务中节点 $i$ 的数据发送能耗为 $e_i(k)$，则跨任务能量状态满足
\begin{equation}
  E_i(k+1)=E_i(k)-e_i(k),\qquad E_i(1)=E_i^{(0)}.
  \label{rq3:eq:energy-evolution}
\end{equation}
任务内每执行一次发送便立即扣减相应电量，式~\eqref{rq3:eq:energy-evolution} 是任务尺度上的汇总形式。

在时刻 $t$，可选动作包括发送候选 $m$、等待长度 $w$ 以及停止。发送已经尝试过的候选表示重传，发送尚未尝试的候选表示新增证据传输。若当前已执行尝试数为 $n_{k,t}$、尝试数上限为 $A_{\max}$，则发送动作的可行集合为
\begin{equation}
  \begin{split}
  \actions^{\mathrm{send}}_{k,t}=\{m\in\mathcal{M}_k:\;&
  \tau_{k,m}\le D_k-t,\\[-2pt]
  &e_{k,m}^{\mathrm d}\le E_{s(m)}(k,t),\quad n_{k,t}<A_{\max}\}.
  \end{split}
  \label{rq3:eq:action-mask}
\end{equation}
等待动作满足 $0<w\le D_k-t$，停止动作始终可执行。硬约束直接排除无法完成的动作，避免仅通过代价惩罚近似电量和截止期限制。

\subsection{优化目标与求解思路}
\label{rq3:subsec:objective}

记策略 $\pi$ 在任务 $k$ 上的停止时刻为 $T_k$，终止接收集合为 $S_k^{\mathrm{end}}$，该任务总通信能耗为
\begin{equation}
  E_k^{\mathrm{comm}}=\sum_{a\in\mathcal{T}_k}
    \left(e_{k,m_a}^{\mathrm d}+e_{k,m_a}^{\mathrm f}\right),
  \label{rq3:eq:total-communication-energy}
\end{equation}
其中 $\mathcal{T}_k$ 是实际执行的数据发送尝试集合。以正确完成收益与通信消耗之间的权衡为目标，可将有限任务周期内的问题写为
\begin{equation}
  \begin{aligned}
    \max_{\pi}\quad &\frac{1}{N_{\mathrm{task}}}\E_{\pi}\left[
      \sum_{k=1}^{N_{\mathrm{task}}}\left(u_k(S_k^{\mathrm{end}})-\lambda_E
      \frac{E_k^{\mathrm{comm}}}{E_{\mathrm{ref}}}\right)\right] \\
    \text{s.t.}\quad &\sum_{k=1}^{N_{\mathrm{task}}} e_i(k)\le E_i^{(0)},\quad i\in\mathcal{N},\\
    &0\le T_k\le D_k,\quad |\mathcal{T}_k|\le A_{\max},\\
    &a_{k,t}\text{ 仅依赖 }\history_{k,t}\text{，且满足式~\eqref{rq3:eq:action-mask}}.
  \end{aligned}
  \label{rq3:eq:global-problem}
\end{equation}
其中 $E_{\mathrm{ref}}>0$ 为固定能量尺度，$\lambda_E\ge0$ 为系统设计中的能耗偏好。式~\eqref{rq3:eq:global-problem} 描述研究目标，并不意味着能够直接获得其全局最优解：真实任务效用在线不可见，未来数据与反馈具有随机性，节点电量还会将不同任务的决策关联起来。

本章采用分层近似求解。首先利用语义效用网络替代不可观测的即时任务正确性；其次把可见历史压缩为接收信念与资源状态；最后以随资源余量变化的能量价格和有限动作前瞻选择当前动作。这样可以保留式~\eqref{rq3:eq:global-problem} 中的交付收益与资源约束，同时避免枚举完整任务周期的反馈分支。后续算法中的能量权重是局部控制参数，需要通过验证任务选择，不与式~\eqref{rq3:eq:global-problem} 的全局最优乘子等同。

\begin{table}[t]
  \centering\small
  \bicaption{语义交付控制模型的主要符号}{Main notation of the semantic delivery control model}
  \label{rq3:tab:notation}
  \begin{tabularx}{\linewidth}{@{}lY@{}}
    \toprule
    符号 & 含义\\
    \midrule
    $N,\ N_{\mathrm{task}}$ & 感知节点数及连续执行的任务总数\\
    $\mathcal{M}_k,\ K_k$ & 任务候选证据集合及其规模\\
    $S_{k,t},\ b_{k,t}(S)$ & 实际接收集合及其条件概率\\
    $c_k(S),\ u_k(S)$ & 接收集合对应的预测类别与真实任务效用\\
    $\hat q_\theta(S),\ q_k(S)$ & 原始预测效用与一致性修正效用\\
    $\bar r_{k,m},\ p_{k,m}$ & 预测链路速率与单次送达概率\\
    $\tau_{k,m},\ D_k$ & 数据发送时长与任务截止时间\\
    $\bar d_{\mathrm f},\epsilon_{\mathrm f}$ & 反馈平均额外时延与反馈丢失概率\\
    $E_i(k,t),\ Z_i(k)$ & 节点剩余发送能量与虚拟能量队列\\
    $\pi_{k,m}$ & 一次发送及其反馈的归一化资源价格\\
    $\eta_E,\eta_T$ & 局部控制的能量权重与时间权重\\
    \bottomrule
  \end{tabularx}
\end{table}
